% Generated draft: requires mathematical and source review.
\documentclass[11pt]{article}
\usepackage[margin=25mm]{geometry}
\usepackage{fontspec}
\setmainfont{Latin Modern Roman}
\usepackage{amsmath,amssymb,mathtools,unicode-math}
\setmathfont{Latin Modern Math}
\usepackage{enumitem,needspace,xurl,hyperref}
\setlength{\parindent}{0pt}
\setlength{\parskip}{5pt}
\setlength{\emergencystretch}{3em}
\newcommand{\N}{\mathbb{N}}
\newcommand{\Z}{\mathbb{Z}}
\newcommand{\Q}{\mathbb{Q}}
\newcommand{\R}{\mathbb{R}}
\newcommand{\C}{\mathbb{C}}
\newcommand{\F}{\mathbb{F}}
\newcommand{\A}{\mathbb{A}}
\newcommand{\PP}{\mathbb{P}}
\newcommand{\E}{\mathbb{E}}
\newcommand{\eps}{\varepsilon}
\newcommand{\dd}{\,\mathrm{d}}
\begin{document}
\section*{UnsolvedMath ID: 20003102}
\textbf{Source problem code:} catalogue entry
\subsection{Definitions and mathematical statement}
Let \(X,Y\) be marked hyperbolic structures on a fixed finite-type surface. Write \(\Lambda(X,Y)\) for the canonical maximally stretched lamination for the directed Thurston problem from \(X\) to \(Y\), as distinguished from a complete lamination chosen to define an individual stretch segment. Let \(q=q(X,Y)\) be the initial quadratic differential of the Teichm\"uller geodesic, with horizontal and vertical measured foliations \(\mathcal H(q),\mathcal V(q)\). The source asks for useful comparisons between these laminations/foliations, including the stretch-segment and high-energy harmonic-map viewpoints, rather than asserting that their supports agree. Foliation supports are compared after geodesic straightening in a chosen hyperbolic metric.

The supplied aid treats a fixed closed surface \(X\) of genus at least two and variable \(Y\ne X\). It combines Thurston's usually-simple theorem with Teichm\"uller polar coordinates to report that, for almost every endpoint \(Y\) in the relevant Lebesgue measure class, \(\Lambda(X,Y)\) is a simple closed curve while neither \(\mathcal H(q)\) nor \(\mathcal V(q)\) has simple-closed-curve support. Thus literal support equality with either foliation is generically false in this setting. The full measure convention and usually-simple theorem hypotheses are not reproduced in the record, and no analogous all-punctured-surface statement is added.

The aid also gives a variational comparison. For nonzero measured foliations \(F\), interpreted as measured laminations for hyperbolic length, set
\[
M_{\ell}(X,Y)=\sup_{[F]}\frac{\ell_Y(F)}{\ell_X(F)}=e^{d_{\rm Th}(X,Y)},\qquad
M_{\rm Ext}(X,Y)=\sup_{[F]}\frac{\operatorname{Ext}_Y(F)}{\operatorname{Ext}_X(F)}=e^{2d_{\rm T}(X,Y)}.
\]
Use the horizontal/vertical convention in which \(\mathcal H(q)\) is the extremal-length-ratio maximizer. The reported nonnegative argmax gap can be written
\[
\Delta(X,Y)=\log M_\ell(X,Y)
-\log\frac{\ell_Y(\mathcal H(q))}{\ell_X(\mathcal H(q))}\ge0.
\]
It is scale-independent in \(\mathcal H(q)\), and vanishes exactly when this extremal-length maximizer is also a hyperbolic-length-ratio maximizer. The aid reports \(\Delta(X,Y)>0\) for almost every \(Y\) in its closed-surface setting. Vanishing tests common variational maximization, not automatically equality with every leaf of the canonical maximal lamination or a chosen complete stretch lamination. The diagonal \(X=Y\) requires a separate convention, since its initial nonzero Teichm\"uller direction is not determined.

The source mentions Minsky's approximate comparison of high-energy harmonic maps, their Hopf foliations, Teichm\"uller maps and stretch maps. No precise estimates or full theorem hypotheses are supplied. The generic mismatch and gap therefore answer only particular exact-equality questions, leaving quantitative, coarse or asymptotic comparisons of the original kinds open in the supplied account.
\subsection{Short English statement}
Compare maximal hyperbolic stretch data with the Teichm\"uller geodesic foliations. The supplied closed-surface result rules out generic support equality and gives an explicit gap testing whether an extremal-length maximizer also maximizes hyperbolic length ratio; broader harmonic-map comparisons remain unspecified.
\subsection{Sources}
\begin{itemize}
\item K. Rafi, Problem 2.12, AIM workshop \emph{Lipschitz metric on Teichmueller space}, \url{https://aimath.org/pastworkshops/lipschitzteichproblems.pdf}. The source names Minsky's high-energy harmonic-map comparisons without a full reference.
\item The research summary invokes Thurston's usually-simple theorem and Teichm\"uller polar coordinates, without full bibliographic entries or a precise measure-class theorem statement.
\item Local UnsolvedMath registry, permanent ID 20003102, exact code \texttt{catalogue entry}; canonical location \texttt{aim-topology-notes.json}, notes[189], dated 14 August 2026; statement labeled exact.
\item Research attempt 1: partial result, problem partially solved at the run, rigor label proved, novelty confidence low; registry status partially solved. The generic mismatch and argmax-gap account are machine-generated research aids, not peer reviewed; independent verification is required before citation.
\end{itemize}
\end{document}
